% Pollis New Methods: arcing classifiers (arc-fs and cost-complexity pruning) — draft 1 (11 October 2026)
% Compile with: latexmk -pdf nm_arcing.tex && latexmk -c
% Shared look for the Pollis New Methods documents (same fonts and colours as the brochure and the workflow handouts).
\documentclass[10pt,a4paper]{article}

\usepackage[a4paper,left=18mm,right=18mm,top=15mm,bottom=17mm,footskip=8mm]{geometry}
\usepackage[T1]{fontenc}
\usepackage[default,semibold]{sourcesanspro}
\usepackage{textcomp}
\usepackage[italic]{mathastext}
\usepackage{amsmath}
\usepackage{microtype}
\usepackage[table]{xcolor}
\usepackage{tikz}
\usepackage[most]{tcolorbox}
\usepackage{fontawesome5}
\usepackage{tabularx,booktabs}
\usepackage{enumitem}
\usepackage{fancyhdr}
\usepackage[hidelinks]{hyperref}

% ---------------------------------------------------------------- colours (as in the brochure)
\definecolor{navy}{HTML}{1C3A5E}
\definecolor{pblue}{HTML}{3B78B5}
\definecolor{amber}{HTML}{F2B33D}
\definecolor{deepamber}{HTML}{D9902A}
\definecolor{pale}{HTML}{EEF3F9}
\definecolor{ink}{HTML}{2A2F36}
\definecolor{logonavy}{HTML}{173D63}
\definecolor{logodark}{HTML}{C98A1E}
\definecolor{logoamber}{HTML}{F2B446}
\definecolor{logocream}{HTML}{FEEFD2}
\definecolor{soft}{HTML}{5E6773}
\definecolor{agree}{HTML}{2F7D4F}
\definecolor{differ}{HTML}{B23B3B}
\color{ink}

\setlength{\parindent}{0pt}
\setlength{\parskip}{4pt}
\setlist{nosep,leftmargin=*,itemsep=2pt}
\renewcommand{\arraystretch}{1.15}

\pagestyle{fancy}
\fancyhf{}
\renewcommand{\headrulewidth}{0pt}
\fancyfoot[L]{\footnotesize\color{soft}\textbf{Pollis}\enspace New Methods\enspace·\enspace\nmtheme}
\fancyfoot[R]{\footnotesize\color{soft}\thepage}
\newcommand{\nmtheme}{}
\newcommand{\nmdate}{11 October 2026}

% The Pollis logo, redrawn so it prints sharply at any size
\newcommand{\pollislogo}[1]{%
  \begin{tikzpicture}[x=#1/512,y=-#1/512,baseline=(current bounding box.center)]
    \fill[logonavy] (0,0) rectangle (512,512);
    \fill[logodark] (36,422) -- (104,256) -- (172,422) -- cycle;
    \fill[logoamber] (104,422) -- (203,180) -- (302,422) -- cycle;
    \fill[logocream] (203,422) -- (339,89) -- (476,422) -- cycle;
  \end{tikzpicture}}

% Top of the first page: logo, name; then the kicker, title and a one-line promise
\newcommand{\nmhead}[3]{%
  \renewcommand{\nmtheme}{#1}%
  \thispagestyle{fancy}%
  \begin{minipage}[c]{13mm}\pollislogo{11mm}\end{minipage}%
  \begin{minipage}[c]{\dimexpr\linewidth-13mm\relax}
    {\color{navy}\Large\bfseries Pollis}\enspace{\color{soft}New Methods}\hfill
    {\color{deepamber}\small\bfseries\MakeUppercase{#1}}
  \end{minipage}\par\vspace{7mm}
  {\color{navy}\fontsize{22}{26}\selectfont\bfseries #2}\par\vspace{4pt}
  {\color{soft}\large #3}\par\vspace{5mm}}

\newcommand{\nmsection}[1]{\par\vspace{5pt}{\color{navy}\large\bfseries #1}\par\vspace{1pt}}
\newcommand{\nmsub}[1]{\par\vspace{2pt}{\color{deepamber}\bfseries #1}\par}
\newcommand{\gt}{\,{\color{deepamber}\bfseries ›}\,}
\newcommand{\menupath}[1]{{\bfseries\color{navy}#1}}
\newcommand{\fn}[1]{\texttt{\small #1}}
% Citations: \nmcite{key}{Breiman (2001)} links to the entry \nmref{key}{...} in the References
\newcommand{\nmcite}[2]{\hyperlink{ref:#1}{\color{pblue}#2}}
\newcommand{\nmref}[2]{\raisebox{1.1em}[0pt][0pt]{\hypertarget{ref:#1}{}}#2\par}
\newcommand{\nmdoi}[1]{\href{https://doi.org/#1}{\color{pblue}doi:#1}}
% Where a notebook is in Pollis: its title in Browse Methods, and its file
\newcommand{\browse}[2]{\menupath{Compose \gt New Methods \gt Browse Methods}: \emph{#1} (\fn{#2.ipynb})}

% Key facts at the top: source, missing in Julia, notebook, result
\newcommand{\factsheet}[4]{%
  \begin{tcolorbox}[enhanced, colback=pale, colframe=pale, boxrule=0pt, sharp corners,
    left=7pt, right=7pt, top=5pt, bottom=5pt, fontupper=\small]
  \begin{tabularx}{\linewidth}{@{}>{\leavevmode\bfseries\color{navy}}p{27mm}>{\leavevmode}X@{}}
  Tested against & #1\\
  Missing in Julia & #2\\
  Notebook & #3\\
  Result & #4
  \end{tabularx}
  \end{tcolorbox}}

% Boxes
\newtcolorbox{pbox}[2][]{enhanced, colback=pale, colframe=amber, boxrule=0pt,
  leftrule=3pt, sharp corners, left=7pt, right=7pt, top=6pt, bottom=6pt,
  fontupper=\small\raggedright, before upper={{\color{navy}\normalsize\bfseries #2}\par\smallskip}, #1}
\newtcolorbox{nbox}[2][]{enhanced, colback=navy, colframe=navy, boxrule=0pt,
  sharp corners, left=8pt, right=8pt, top=7pt, bottom=7pt, coltext=white,
  fontupper=\small\raggedright, before upper={{\color{amber}\normalsize\bfseries #2}\par\smallskip}, #1}

% Value tables: published beside Julia
\newenvironment{valuetable}[1]{%
  \par\small\tabularx{\linewidth}{#1}\toprule}{\bottomrule\endtabularx\par}
\newcommand{\ok}{{\color{agree}\faCheck}}
\newcommand{\no}{{\color{differ}\faTimes}}

% End of the document: credit and licence
\newcommand{\nmcredit}{%
  \par\vfill
  \begin{minipage}[b]{0.78\linewidth}\footnotesize\color{soft}
  Methods written and tested with Claude Code (Anthropic) in Pollis, from the published papers and their data.\par
  \textcopyright\ 2026 Antonio Saragga Seabra\enspace·\enspace Licensed under \href{https://creativecommons.org/licenses/by-sa/4.0/}{CC BY-SA 4.0}\enspace·\enspace Draft, \nmdate
  \end{minipage}\hfill
  \begin{minipage}[b]{0.18\linewidth}\raggedleft\pollislogo{9mm}\par\vspace{2pt}{\footnotesize contact@pollis.org}\end{minipage}}


\begin{document}

\nmhead{Boosting}
  {arc-fs: adaptive resampling with pruned trees}
  {The AdaBoost of Breiman's comparisons, tested against his published test errors on 20 data sets.}

\factsheet
  {\nmcite{breiman2001}{Breiman (2001)}, \emph{Random forests}, Table 2, the AdaBoost column (p.\,11) and Section 6 (p.\,15); \nmcite{breiman1998}{Breiman (1998)},
   \emph{Arcing classifiers}, Tables 3 and 5. 15 real data sets from PollisDatasets, four simulated as Breiman did, and the
   zip-code digits of \nmcite{hastie2009}{Hastie, Tibshirani and Friedman (2009)} (downloaded only if the user agrees).}
  {arc-fs, and CART pruning by cost complexity on a test sample. DecisionTree.jl grows the trees only.}
  {\browse{arc-fs: adaptive resampling with pruned trees}{arcfs}}
  {18 of 20 test errors within simulation noise of the published ones.}

\nmsection{The problem}
arc-fs is the version of AdaBoost of \nmcite{freund1996}{Freund and Schapire (1996)} behind Breiman's comparisons: the ``AdaBoost'' column against
which random forests were first judged. Each step draws a training set that favours the cases earlier trees got wrong, grows a
CART tree, prunes it, and the trees vote with weights. No Julia package has arc-fs, and none prunes a CART tree by cost
complexity on a test sample. Without both, the classic comparison of boosting and forests cannot be rerun in Julia.

\nmsection{The method}
\nmsub{Cost-complexity pruning}
CART, \nmcite{breiman1984}{Breiman, Friedman, Olshen and Stone (1984)}, Ch.\,3, repeatedly collapses the weakest link, the node $t$ with the smallest
$g(t) = (R(t) - R(T_t))/(|T_t| - 1)$, where $R(t)$ counts the cases misclassified at $t$ as a leaf, $R(T_t)$ those
misclassified by the branch below it and $|T_t|$ its leaves. \fn{prune} keeps, from this sequence, the subtree with fewest
errors on a second sample, the smallest on ties.
\nmsub{arc-fs}
A probability $p(n)$ sits on each training case, equal at the start. At step $k$, draw a training set with these probabilities,
grow a tree and prune it on a second draw. With $d(n) = 1$ for the cases the tree misclassifies,
\[ \varepsilon_k = \textstyle\sum p(n)\,d(n), \qquad \beta_k = (1-\varepsilon_k)/\varepsilon_k, \qquad
   p(n) \leftarrow p(n)\,\beta_k^{d(n)}, \text{ normalised}, \]
and the trees vote with weights $\log\beta_k$. When $\varepsilon_k$ reaches $\tfrac12$, or is 0, the probabilities are reset
and the construction restarts; that step's tree does not vote.

\nmsection{What the tests show}
On each small data set a random 10\% is the test set and 50 trees are combined; 100 runs are averaged. Letters, sat-images and
zip-code have their own test sets and are one run each, with a binomial standard error. A value is \emph{within noise} when
the two differ by less than two standard errors of a difference ($\sqrt2$ s.e.).
\begin{valuetable}{@{}lXl@{}}
\textbf{Data sets} & \textbf{What is compared} & \textbf{Within noise}\\\midrule
13 small real sets & average test error of 100 runs & 13 of 13\\
letters, sat-images, zip-code & one run on their own test sets & 3 of 3\\
twonorm, threenorm, ringnorm, waveform & average test error of 50 runs & 2 of 4\\
\midrule
All & 20 test errors & 18 of 20\\
\end{valuetable}

\nmsub{Where they differ}
Twonorm 4.6 vs 4.9 and ringnorm 6.2 vs 6.9: below Breiman by 2.8 and 4.4 standard errors. Among those that agree, breast
cancer (3.7 vs 3.2) and diabetes (25.3 vs 26.6) came closest to the limit. The notebook's example shows the point of arcing:
on one split of the breast cancer data, the vote of a few dozen pruned trees halves the error of a single one.

\nmsection{Still open}
\begin{itemize}
  \item \textbf{Twonorm and ringnorm.} The random forests of the Out-of-bag notebook are below Breiman on the same two sets.
    A gap shared by two methods points to how the trees are grown, or to the generators, rather than to arc-fs. Not explained.
  \item \textbf{Details the papers leave open:} whether a restarting step's tree joins the vote (here it does not), and how
    CART chose among equally good pruned subtrees (here the smallest).
  \item \textbf{The data.} Soybean has 683 cases, not 685, and its categorical inputs are split as ordered codes, since
    DecisionTree.jl has no splits on subsets of categories. Missing values are replaced by medians; Breiman does not say
    what he did.
\end{itemize}

\nmsection{Further considerations}
\begin{itemize}
  \item \textbf{Pruning on its own.} Cost-complexity pruning is useful well beyond arc-fs: it is how a single CART tree is
    usually chosen, and it is missing from DecisionTree.jl. It could be offered upstream with these tests.
  \item \textbf{Categorical splits.} Splits on subsets of categories, which Breiman's CART had, would remove the soybean
    caveat and are a natural next contribution.
  \item \textbf{Next.} The arcing paper's arc-x4 and bagging columns, from the same tables, to rerun the whole comparison.
\end{itemize}

\nmsection{References}
{\footnotesize
\nmref{breiman1998}{Breiman, L. (1998). Arcing classifiers. \emph{The Annals of Statistics} 26(3), 801--849. \nmdoi{10.1214/aos/1024691079}}
\nmref{breiman2001}{Breiman, L. (2001). Random forests. \emph{Machine Learning} 45(1), 5--32. \nmdoi{10.1023/A:1010933404324}}
\nmref{breiman1984}{Breiman, L., Friedman, J.\,H., Olshen, R.\,A. and Stone, C.\,J. (1984). \emph{Classification and Regression Trees}. Wadsworth.}
\nmref{freund1996}{Freund, Y. and Schapire, R.\,E. (1996). Experiments with a new boosting algorithm. \emph{Proceedings of the Thirteenth International Conference on Machine Learning}, 148--156.}
\nmref{hastie2009}{Hastie, T., Tibshirani, R. and Friedman, J. (2009). \emph{The Elements of Statistical Learning}, 2nd ed. Springer: the zip-code data.}
}

\nmcredit
\end{document}
